\section*{Chapter \ref{ch:ll:build-tick-to-trade-measured} --- Build: Tick-to-Trade, Measured}

\begin{solution}{exo:ll:build-tick-to-trade-measured:1}
Book, strategy, risk and gateway take $160 + 100 + 47 + 33 \approx \qty{340}{\nano\second}$ at their medians, under a quarter of
the \qty{1.5}{\micro\second} that the in-process stages add up to (decode and ring included), and under a tenth of the in-process
median over all orders, which includes the waits of second orders.
\end{solution}

\begin{solution}{exo:ll:build-tick-to-trade-measured:2}
A stage is the difference of two readings taken on different threads, often on different cores. The \omterm{def:ll:cpu-microarchitecture:tsc}{time-stamp counter} is one
counter, invariant and synchronised across the cores of this machine (chapter~5), so the difference is a duration. Per-thread
clocks with different offsets or rates would make stages drift, and some would come out negative.
\end{solution}

\begin{solution}{exo:ll:build-tick-to-trade-measured:3}
The $k$-th event of the packet is published after $k \times \qty{100}{\nano\second}$, so the five see 100 to 500, on average
\qty{300}{\nano\second}. On the engine side the same five events are taken one after another, each waiting for the handling of
the ones before it, and for their orders' sends.
\end{solution}

\begin{solution}{exo:ll:build-tick-to-trade-measured:4}
The risk stage, measured from the decision to the check of the second order: it would include the whole send of the first
order, about \qty{4}{\micro\second}. With half of the orders in that case, the stage's median would land between the two groups,
anywhere from the check's own tens of nanoseconds to the send's microseconds: the attribution would be wrong, not the gate.
\end{solution}

\begin{solution}{exo:ll:build-tick-to-trade-measured:5}
Keeping about \qty{100}{\nano\second} of decoding proper: $0.9 + 0.8 + (0.10 + 0.16 + 0.10 + 0.05 + 0.03) \approx
\qty{2.1}{\micro\second}$, within the budget, with about \qty{0.9}{\micro\second} to spare.
\end{solution}

\begin{solution}{exo:ll:build-tick-to-trade-measured:6}
Both stamps are taken by software around system calls: the exchange's before \texttt{sendto}, the feed thread's after
\texttt{recv}, so the stage contains the sender's call, the kernel's loopback path and the receiver's call, with no reading in
between. A hardware time stamp taken by the card when the packet is on the wire splits it into the sender's part and the
receiver's part, and on a real network adds the wire itself.
\end{solution}

\begin{solution}{exo:ll:build-tick-to-trade-measured:7}
One ring message per packet with up to five events; the engine runs them in order, so the orders are unchanged (the test's hash
checks it). The decode stage of all but the last event disappears, since the batch is published at once, and the ring stage
loses the per-message overhead; the queue behind earlier orders' sends remains.
\end{solution}

\begin{solution}{exo:ll:build-tick-to-trade-measured:8}
It measures only the in-process part, leaving out the network card, the kernel's receive and transmit (on this laptop, most of
the total) and the wire; it stops at the call to \texttt{send}, not when the order leaves; and an average says nothing about the
tail, where races are lost. Report the percentiles of the interval from the packet on the wire to the order on the wire.
\end{solution}

\begin{solution}{pb:ll:build-tick-to-trade-measured:1}
\begin{enumerate}
\item Receive (exch to recv), decode (recv to pub), ring (pub to pop), book (pop to book), strategy (book to decide), risk (decide
to risk), gateway (risk to encode), transmit (encode to venue).
\item Tick to trade: about \qty{8.8}{\micro\second}, \qty{79}{\micro\second} and \qty{3.5}{\milli\second}; in process: about
\qty{3.7}{\micro\second}, \qty{42}{\micro\second} and \qty{3.5}{\milli\second}.
\item Transmit has the largest median (about \qty{2.9}{\micro\second}); the ring the largest 99.9th percentile (about
\qty{3.5}{\milli\second}).
\item When the engine thread's CPU is taken away (the kernel, the hypervisor, another task), events wait on the ring until the
thread runs again: the machine's scheduling, not the code.
\item The strategy and the \omterm{def:ll:pre-trade-risk-and-kill-switches:gate}{risk gate} meet their targets; decode misses by a factor of about 7, book by a third, receive by 1.6 and
transmit by 3.6.
\item The stages of one order compose exactly into its own total; using the first orders keeps the composition consistent with
the stage attribution, while the tick-to-trade distribution includes the second orders' waits.
\item Quantiles of a sum are not sums of quantiles (chapter~1): the slow cases of different stages rarely coincide, and when
they do, they are not the same messages.
\item That the path meets, or does not meet, a stated performance target on stated hardware, stage by stage: a property of the
system that no functional test sees.
\item $0.9 + 0.8 + 1.5 \approx \qty{3.2}{\micro\second}$ at the median.
\item About two-thirds of the \qty{1.5}{\micro\second} in process is queueing (the packet's events waiting for each other in
decode and on the ring); the rest is computation, of which the book's components take about \qty{340}{\nano\second}.
\item They would remove the preemption that makes the ring's tail; they would not change the medians of the computation, nor the
kernel's receive and transmit.
\item It adds figures from a budget to measured ones; only a run of this harness on the target server, with its card and its
tuning, measures it.
\item \textbf{Named result.} On this laptop, otherwise idle, over 10\,500 orders: receive about 1.4, decode 0.44, ring 0.76, book
0.16, strategy 0.10, risk 0.05, gateway 0.03 and transmit \qty{2.9}{\micro\second} at the median; tick to trade about
\qty{8.8}{\micro\second} at the median, \qty{79}{\micro\second} at the 99th and \qty{3.5}{\milli\second} at the 99.9th percentile;
projected for a tuned server with \omterm{def:ll:linux-tuning:polling}{kernel bypass}, about \qty{3.2}{\micro\second} at the median, \qty{2.1}{\micro\second} without
the packet's queue.
\item About \qty{340}{\nano\second} of computation by the book's components against some \qty{1.2}{\micro\second} of decoding and
waiting in the \qty{1.5}{\micro\second} of first-order stages; the in-process median of all orders, \qty{3.7}{\micro\second}, is
mostly waiting.
\item Transmit and receive: \omterm{def:ll:linux-tuning:polling}{kernel bypass} (chapter~13 and One Quant Book~14); the ring's tail: \omterm{def:ll:linux-tuning:isolation}{core isolation} (chapter~13); the
queue: batching the hand-over (chapters~12 and 18).
\item Otherwise the timing could be of a path that computes something else, or that skips work when something goes wrong.
\item An allocation on the path is a latency hazard (chapter~6) that a change can introduce silently; the count catches it on
the next run.
\item The receive and transmit stages with hardware time stamps, and the tails with the path's CPUs isolated.
\item The per-component figures of the earlier chapters: each was measured alone, with warm caches and no neighbours.
\item Measure it on the path that computes the right answer, stage by stage, with one clock, in percentiles.
\end{enumerate}
\end{solution}

\begin{solution}{iq:ll:build-tick-to-trade-measured:1}
The time from the market event reaching the firm (ideally, the packet on the wire at the network card) to the resulting order
leaving it (the order on the wire). Start and stop with hardware time stamps where possible; otherwise say exactly which
software points are used and what they leave out.
\iqlookfor{wire to wire, and honesty about the points.}
\end{solution}

\begin{solution}{iq:ll:build-tick-to-trade-measured:2}
Read one shared, cheap clock (the \omterm{def:ll:cpu-microarchitecture:tsc}{time-stamp counter}) at a few stage boundaries, carry the readings with the message in
preallocated records, write them out off the hot path (a \omterm{def:ll:logging-capture-and-replay:binlog}{binary log}), and join them by message identity afterwards; attribute
stages per message and report percentiles.
\iqlookfor{few points, one clock, off-path recording, per-message attribution.}
\end{solution}

\begin{solution}{iq:ll:build-tick-to-trade-measured:3}
At what the path waits for rather than what it computes: preemption and interrupts on its CPUs (isolation), \omterm{def:ll:cpp-for-latency-i-memory:fault}{page faults} and
allocation, queues behind bursts, the kernel's network stack, garbage collection or logging on the path. The stage whose tail
grows tells which.
\iqlookfor{scheduling and queueing before code.}
\end{solution}

\begin{solution}{iq:ll:build-tick-to-trade-measured:4}
Wire to card and DMA, then a kernel-bypass receive (under a microsecond); decode tens of nanoseconds; book and strategy tens to a
few hundred; risk tens; encode tens; kernel-bypass send and card to wire (under a microsecond). A few microseconds in all, most of
it in the network path.
\iqlookfor{per-stage orders of magnitude, and where the time goes.}
\end{solution}

\begin{solution}{iq:ll:build-tick-to-trade-measured:5}
From the competitive requirement at the percentiles that matter, allocated to stages with the union bound for the tails
(chapter~1); proved by a harness that stamps every stage on the production path and checks every release against the budget.
\iqlookfor{targets per stage and percentile, and a harness that checks them.}
\end{solution}

\begin{solution}{iq:ll:build-tick-to-trade-measured:6}
Hardware time stamps at the card for wire in and out; the \omterm{def:ll:cpu-microarchitecture:tsc}{time-stamp counter} at each stage boundary, calibrated against a
UTC-traceable clock; stamps carried with the message and written to a \omterm{def:ll:logging-capture-and-replay:binlog}{binary log} by a separate thread; joined by sequence number
and order identifier; per-stage histograms per release, with alerts on the tails; the same harness in the performance gate.
\iqlookfor{clocks, joins, histograms, and the link to the release process.}
\end{solution}
